\documentclass[a4paper,11pt]{ltjsarticle}
\usepackage[margin=25mm,headheight=16pt]{geometry}
\usepackage{amsmath,amssymb,amsthm,mathtools,bm}
\usepackage{booktabs,longtable,array,enumitem}
\usepackage{xcolor}
\usepackage[unicode,colorlinks=true,linkcolor=blue!40!black,citecolor=blue!40!black,urlcolor=blue!45!black]{hyperref}
\usepackage{fancyhdr}
\usepackage{microtype}
\setlength{\parindent}{1em}
\setlength{\parskip}{0.25em}
\setlist{itemsep=0.2em,topsep=0.4em}
\allowdisplaybreaks[2]
\pagestyle{fancy}
\fancyhf{}
\fancyhead[L]{\small 3次超曲面のPoisson構造と非可換化}
\fancyhead[R]{\small 2026-11-08 / v2}
\fancyfoot[C]{\thepage}
\newcommand{\CC}{\mathbb C}
\newcommand{\ZZ}{\mathbb Z}
\newcommand{\PP}{\mathbb P}
\newcommand{\OO}{\mathcal O}
\newcommand{\cK}{\mathcal K}
\newcommand{\cP}{\mathcal P}
\newcommand{\cD}{\mathcal D}
\newcommand{\cC}{\mathcal C}
\newcommand{\cE}{\mathcal E}
\newcommand{\cF}{\mathcal F}
\newcommand{\cM}{\mathcal M}
\newcommand{\fm}{\mathfrak m}
\newcommand{\Perf}{\operatorname{Perf}}
\newcommand{\Ku}{\operatorname{Ku}}
\newcommand{\MF}{\operatorname{MF}}
\newcommand{\HH}{\operatorname{HH}}
\newcommand{\QGr}{\operatorname{QGr}}
\newcommand{\qgr}{\operatorname{qgr}}
\newcommand{\Dsg}{D_{\mathrm{sg}}}
\newcommand{\RHom}{\operatorname{RHom}}
\newcommand{\REnd}{\operatorname{REnd}}
\newcommand{\Hom}{\operatorname{Hom}}
\newcommand{\Ext}{\operatorname{Ext}}
\newcommand{\Ind}{\operatorname{Ind}}
\newcommand{\Def}{\operatorname{Def}}
\newcommand{\thick}{\operatorname{thick}}
\newcommand{\Spec}{\operatorname{Spec}}
\newcommand{\Spf}{\operatorname{Spf}}
\newcommand{\Id}{\operatorname{Id}}
\newcommand{\sym}{\operatorname{sym}}
\newcommand{\Cone}{\operatorname{Cone}}
\newcommand{\rank}{\operatorname{rank}}
\newcommand{\Pic}{\operatorname{Pic}}
\newcommand{\otL}{\mathbin{\otimes^{\mathbf L}}}
\newcommand{\DocDate}{2026年11月8日}
\newcommand{\DocVersion}{統合第2版}
\theoremstyle{plain}
\newtheorem{theorem}{定理}[section]
\newtheorem{proposition}[theorem]{命題}
\newtheorem{lemma}[theorem]{補題}
\newtheorem{corollary}[theorem]{系}
\theoremstyle{definition}
\newtheorem{definition}[theorem]{定義}
\newtheorem{problem}[theorem]{問題}
\theoremstyle{remark}
\newtheorem{remark}[theorem]{注意}
\renewcommand{\proofname}{証明}
\hypersetup{pdfauthor={},pdfsubject={3次超曲面のPoisson構造・非可換量子化・rotation付き残余圏}}

\hypersetup{pdftitle={非可換cubicの可換モデル：研究ノート}}
\title{研究ノート\\\large 非可換cubicの可換モデル：証明監査・適用条件・次の問題}
\author{}
\date{\DocDate・\DocVersion}
\begin{document}
\maketitle
\tableofcontents

\section{問いへの現在の答え}

滑らかな可換cubicから来るrotation付きKuznetsov成分には、polarized categorical
Torelliが適用できる \cite[Theorem 1.7]{LZ}。
任意の非可換な組が可換cubicの像に属するという存在命題は、Torelliの一意性とは別である。
AS-regularな周囲と三次中心正則元による広い定義では、その存在命題には滑らかな反例がある。
周囲にCalabi--Yau性を加えても反例は残る。
現行論文の肯定的な存在定理、反例、および未解決範囲は次の通りである。

\begin{longtable}{p{0.30\linewidth}p{0.60\linewidth}}
\toprule
範囲 & 結論と証明の状態\\
\midrule
\endhead
smooth cubicの形式近傍 &
flat graded変形から相対pairを作り、一意なmarked cubicを得る。本文に証明。\\
\(n\ne4\) の形式近傍 &
rotationを忘れても同じ復元が成立し、rotationのliftのmarked同型類は一意。本文に証明。\\
smooth properな代数的pairの族 &
cubic点のétale近傍で可換モデルとpair同値を作れる。実現点集合はZariski開。
本文にmoduliと持ち上げの証明。\\
一Sklyaninブロックと中心三次式 &
\(h^6\ne-27\)、中心三次式の臨界点が原点だけなら、有限パラメータで明示モデルを得る。
本文に生成性と降下の証明。\\
複数の独立な楕円ブロック &
対応する可換cubicの残余圏を全MF圏へ完全忠実に埋め込める。
高々一つを除く平行移動が有限位数なら全生成性も成立し、任意の孤立中心三次式を許す。
二つ以上の無限位数因子では未解決。\\
位数の異なる明示パラメータ &
\(h_j^6=5\) では平行移動点の位数は5。二ブロックのsixfoldと三ブロックのsevenfoldに適用できる。
接線計算の証明と厳密記号検算を備える。\(h=2\) は無限位数で、良い還元での点数から証明する。\\
Fermat中心部分で \(h^6=-27\) &
標準MF圏はnonproper、\(\qgr B\) のhomological dimensionも無限。本文に証明。\\
AS-regular中心商の滑らかな三次元例 &
Fermat cubicのZhang twistはsmooth properで \(\rho^3\simeq[2]\) も満たすが、
Serre--rotation条件を破り、可換pair模型を持たない。論文に証明。\\
Calabi--Yau中心商の滑らかな四次元例 &
六次元CY ambientのZhang twistでも、\(\dim\HH_0^{\rho}=8\) となる。
可換cubic fourfoldの標準値20と異なり、可換pair模型を持たない。論文に証明。\\
smooth properな連結pair族 &
相対的に \(\rho^3\simeq[2]\) があれば不変Hochschild Euler標数は一定。
標準cubic fourfoldの22とCY反例の10は同じ連結族で結べない。論文に証明。\\
本稿の量子化族の有限パラメータ &
相対saturationの構成、非可換変数との混合項、元の底への降下、
可換点を含む特定の成分の全ファイバーでの存在は未解決。\\
\bottomrule
\end{longtable}

上記の結果は研究草稿の定理として証明を備えるが、公刊・査読済みではない。
形式復元、étale局所復元、一ブロックと多ブロックの比較、有限中心の生成定理、
およびrotationの反例の証明は論文に置き、PI中心幾何による別証明は本ノートに置く。
個々の結果の証明状態と重点検証箇所は研究方針に記す。
圏論の形式化や外部査読は行っていない。
BFK/HR/JLLZの全計算の独立再証明は行っておらず、一次資料の結果に依拠する。
先行研究上の新規性は未確認である。
以下では、既知の基礎定理と本稿で組み立てた証明の境界を明記する。

\section{可換模型の適用範囲と量子化成分の境界}
\label{sec:notes-model-boundary}

\subsection{圏の同値とrotation付き同値を区別する}

右Zhang twistの積を \(a*b=a\tau^{\deg a}(b)\) とすると、
\(\operatorname{qgr}R^\tau\simeq\operatorname{qgr}R\) であっても、
untwist後の次数シフトは幾何自己同型 \(f^*=\tau\) を用いた
\(f_*(-\otimes\OO(1))\) になる。
従って残余圏のunderlying dg圏が同じでも、rotationは \(f_*\rho_X\) に変わる。
\(\operatorname{qgr}\) の同値だけでrotation付き存在問題を解いたと判断できない。

\(\zeta^3=1\)、\(\zeta\ne1\) とする。\(N=n+2\) に対する必要条件は
\[
 S_{\mathcal K}\simeq\rho^{3-N}[N-2].
\]
Fermat cubic threefoldに \(\tau=\operatorname{diag}(\zeta,1,1,1,1)\) を用いると、
この関係が破れる。さらにFermat cubic fourfoldに
\(\tau=\operatorname{diag}(\zeta,\zeta^2,1,1,1,1)\) を用いると、
周囲はCalabi--Yauであり、このSerre関係と \(\rho^3\simeq[2]\) の両方を満たすが、
\(\HH_0\) 上のrotation不変部分の次元が8となる。
標準cubic fourfoldの値20と異なるため、rotation付き可換模型は存在しない。
証明と左右twistの規約は、論文の「可換cubic模型の存在を妨げるrotationの不変量」節に置く。

\subsection{連結な相対saturated族を区別する不変量}

smooth properな相対dg圏と \(\rho^3\simeq[2]\) があるとき、
\[
 \chi_\rho=\sum_i(-1)^i\dim\HH_i^{\rho}
\]
は連結な底上で一定である。
相対Hochschild複体の \((1+\rho+\rho^2)/3\) による直和因子がperfectであることを用いる。
各次数のHochschild群の局所自由性は仮定しない。
四次元の標準cubic pairでは \(\chi_\rho=22\)、上のCY反例では \(\chi_\rho=10\) である。
従って、相対周期性を保つ連結smooth proper族では両者は結ばれない。
詳細は論文の「smooth properな連結族で保存される不変量」小節を参照する。

この反例を本稿のPoisson量子化族のファイバーと同定したわけではない。
有限位数Zhang twistのパラメータを任意に動かしても、三次中心元が残るとは限らない。
したがって研究対象は、単にCY中心商の集合ではなく、
零で可換となる指定した中心量子化族と、その相対残余pairを含む成分である。
\(\chi_\rho=22\) は四次元のその成分で必要な確認値だが、可換模型の存在の十分条件ではない。

\subsection{独立ブロックの全生成性に対する二つの方法}

二つ以上の無限位数ブロックでは、粗化次数付きTHCR加群の外部積による生成条件(G)は未解決である。
これを直接証明する方法に加え、既存の完全忠実比較
\(\Phi:\Ku(Y)\to\mathcal K_B\) に次の判定を適用できる。
\(\mathcal K_B\) がproperで、Serre関手が
\(\rho_B^{3-N}[N-2]\) であり、\(\HH^0(\mathcal K_B)\) に非自明な冪等元がなければ、
\(\Phi\) は同値である。
sourceのsmooth proper性から随伴を得て、Serre互換性で直交分解を導き、
\(\HH^0\) の連結性で補圏を消す。
論文の「完全忠実比較を同値へ強める条件」節に証明を置く。

この判定の三条件を全ての無限位数パラメータで検証したわけではない。
また一般のCY中心商には \(\Phi\) 自体の存在を仮定できない。
この区別が、上のCY反例と独立ブロックの条件付き定理を両立させる。
ここで \(\HH^0\) はcohomologyであり、反例に使うhomologyの \(\HH_0\) とは異なる。

\subsection{量子化成分での具体的な着手}

第一に、本稿の中心量子化族について、相対Beilinson仮定と対角の有限証明書を検証し、
相対smooth properな残余dg圏、rotation、および \(\rho^3\simeq[2]\) を同じ底上に構成する。
これにより可換点のétale近傍での存在定理と、連結成分上の \(\chi_\rho\) の制約を同時に使える。
各Artin切断での結果だけを、この代数的相対構成の代わりにしない。

第二に、独立な二つの無限位数ブロックでは、条件(G)の係数付きTHCRによる証明と、
properness・Serre公式・\(\HH^0\) の連結性の検証を比較する。
後者の三条件を示せば、全ての加群を外部積から直接構成する必要はなくなる。
どちらも細分不能な粗化次数の対象と軸torsionを対象範囲から除かない。

第三に、非可換生成元と中心変数の混合項では、相対THCRと相対windowを作り、
translationの消去を底変更と両立させる。
étale局所モデルの構成後にも、同値のmonodromy、元の底への降下、
実現開部分の外のファイバーはそれぞれ別の問題として残る。

\section{量子化定理から復元論へ渡す条件}

論文の量子化定理（第I部）は、零を含むZariski開集合 \(U\) 上で
\[
 B_\lambda=A_\lambda/(\Omega_\lambda),\qquad
 B_0=\CC[x_1,\ldots,x_N]/(F)
\]
を構成し、周囲のKoszul AS-regular性、三次の中心元の正則性、各次数のflatness、
\[
 H_{B_\lambda}(t)=\frac{1-t^3}{(1-t)^N}
\]
を主張する。\(N=n+2\) である。
量子化定理は、\(\qgr B_\lambda\) の圏論的滑らかさを別問題としている。

形式復元に使う入力は、\(\hbar=0\) での完備化の各Artin切断が、
滑らかな \(X=V(F)\) の座標環の次数ごとに有限自由な変形であることだけである。
従って、Poisson分類や中心元構成を復元論の中で再利用しない。
量子化の証明の検証範囲は研究方針に記す。

なお \(\Omega\) の中心持ち上げやパラメータの正規化の選び方によって、
得られる可換モデルの表示は変わり得る。
明示例で使う \(G_h\) は現行論文で固定した中心持ち上げである。

\section{相対残余圏の構成における監査点}

\subsection{bounded derivedとperfectを分ける}

\(R\) が非正則なArtin環のとき、可換な \(X_R\) についてさえ
\(D^b(\operatorname{coh}X_R)\) 全体は \(R\)-smooth properな変形の適切な対象ではない。
用いるのは
\[
 \cP_R=D(\QGr B_R)^c
\]
である。有限パラメータの特異圏との比較では、体上の
\(D^b(\qgr B)\) のOrlov成分を使う。この二つの使い分けが必要である。

\subsection{De Deken--Lowenの適用}

\(\OO_X(i)\) は \(\ZZ\)-generatingで、\(i\le j\) について
\[
 \Ext_X^{1,2}(\OO_X(i),\OO_X(j))=0.
\]
\cite[Theorem 4.5]{DDL} によりgraded algebra deformationからflat abelian deformationを得る。
ここで適用しているのは全 \(\ZZ\)-列のdirected Hom-categoryであり、
短い例外列が \(D^b(X)\) 全体を生成すると仮定していない。

生成子には
\[
 E_R=\bigoplus_{i=0}^{n+1}\OO_R(i)
\]
を使う。特殊ファイバーでBeilinsonの生成子の制限になり、
\cite[Propositions 5.8--5.9]{DDL} でcompact generationを持ち上げる。
Proposition 5.5の導来Homの底変更を、dg reduction functorから構成することで
\[
 \REnd(E_R)\otL_Rk\simeq\REnd_X(E_0)
\]
を合成・単位を保つ代数の比較として得る。
cohomologyの各次元だけの一致では不十分である。

\subsection{nilpotent descentは通常のNakayamaではない}

局所Artin底では \(R\in\thick_R(k)\) である。
\(M\otL_Rk\) のcompactnessを、\(k\otL_R A\) 上で確認すれば
\(M\) の \(A\)-compactnessが従う \cite{Mathew}。
これは \(k\) が \(R\)-perfectであることを必要としない。
同じ補題を、生成子の自己Homとその対角bimoduleの両方へ適用する。
前者はproperness、後者はsmoothnessを与える。

残余圏の生成子は \(p_RE_R\)、rotationは \(\mathbf L_{\OO_R}(1)\) である。
rotationが右直交を保つ計算は \(1\le i\le n-2\) の範囲で行う。
完全忠実性はHom比較のconeの還元、本質的全射性は生成子の像の還元から確認する。

\section{形式復元における監査点}

\subsection{Hochschild cohomologyの低次数}

BFKの元のKoszul公式 \cite[Theorem 5.39]{BFK} を使うと、
\(N=n+2\)、\(n\ge3\) で
\[
 \HH^{<0}=0,\quad \HH^0=\CC,\quad \HH^1=0,
 \qquad
 \HH^2=
 \begin{cases}J(F)_3&n\ne4,\\J(F)_3\oplus\CC^2&n=4.\end{cases}
\]
恒等sectorは \(\HH^{2r}\) に \(J(F)_{3r}\) を与える。
二つの非自明 \(\mu_3\)-sectorは、\(3\mid N\) の場合にのみ
cohomological degree \(N/3\) に現れる。

\begin{center}
\begin{tabular}{ccccc}
\toprule
\(n\)&\(N\)&\(\dim J_3\)&\(\dim\HH^2\)&余分なsectorの位置\\
\midrule
3&5&10&10&なし\\
4&6&20&22&\(\HH^2\) に二方向\\
5&7&35&35&なし\\
6&8&56&56&なし\\
7&9&84&84&\(\HH^3\) に二方向\\
\bottomrule
\end{tabular}
\end{center}

四次元の余分な二方向ではrotationの固有値は \(\zeta,\zeta^2\) である。
HuybrechtsのMukai格子上の作用 \cite[Proposition 3.12]{Huy} と
invariantなCalabi--Yau volumeによる \(\HH^2\simeq\HH_0\) を使う。
従って全次元で \(\HH^2{}^\rho=J(F)_3\) となる。

\subsection{接写像の同型は次元だけではない}

\(F+\varepsilon G\) の接ベクトルが恒等sectorの \([G]\) に写ることが必要である。
\(n\ge4\) では \cite[Lemma 2.19]{HR} が、matrix factorizationの周期性に由来する
kernel morphismとKodaira--Spencer morphismを比較する。
同論文の条件 \(3\le(n+2)/2\) を忘れない。
\(n=3\) は \cite[Theorem 4.4(2)]{JLLZ} で補う。
JLLZの残余圏の例外列の規約は、全体への直線束テンソルで本稿と同一視できる。

\subsection{curvature問題と比較写像}

naiveなMorita変形をHochschild cohomology全体と同一視すると、
curvatureに関する問題がある \cite{KL}。
本稿では、tame compactly generatedな \(\cC=\Ind\Ku(X)\) に対する
\[
 \theta_{\cC}:\operatorname{CatDef}_{\cC}^c
 \longrightarrow\operatorname{CatDef}_{\cC}^{\wedge}
\]
の \((-1)\)-truncated性だけを使う \cite{Lurie,BKP}。
右側で同値となった二つのactual deformationを、左側でも同値と判断できる。
右側の任意の点がactual categoryになるとは仮定していない。

一つの可逆なrotationを持つ変形は、markingを動かす \(\ZZ\)-作用のhomotopy fixed pointである。
その接複体は
\[
 R\Gamma(\ZZ,CC^\bullet(\Ku(X))[2])
\]
であり、\(\HH^1=0\) によって一次変形は \(\HH^2{}^\rho\) と一致する。
この定式化は、\(\rho\) を厳密な有限群作用とみなすことを要求しない。

\subsection{接空間から全Artin次数へ進む理由}

small extension \(R'\to R\) の核を \(I\) とする。
幾何側で一つliftを作り、既に与えられたpairのliftと、還元の同定を固定した上で比較する。
formal moduli problemのhomotopy fibre内で、存在するliftの差は
\(H^0(T)\otimes I\) のtorsorに属する。
接写像の同型で幾何側を修正し、この差を消す。
これが全射性のArtin帰納法である。

単射性には、還元上に現れたpairの自己同型を持ち上げる操作が必要である。
\(\HH^1=0\) により恒等に還元されるinvertible bimoduleは対角と同型となり、
\(\HH^0_R=R\) により残るpair自己同型はscalar unitになる。
markingを固定したscalarは \(1+\fm_R\) に属し、\(R'\) へ持ち上がる。

sevenfoldの \(\HH^3=\CC^2\) はこの証明を妨げない。
既に存在する二つのliftを比較しているので、\(\HH^3\) 全体の消滅もformalityも不要である。
ただし高次自己同型は残るため、full higher groupoidの同値という結論にはしない。

\section{étale局所復元の入力条件と証明の境界}

\subsection{圏そのもののmoduliを構成しない}

smooth properな代数的族 \(\cK_S\) と、方程式空間 \(U\) 上の普遍cubicの残余圏を
先に固定する。そのtensor product
\[
 \cK_S^{\mathrm{op}}\otL\Ku(\mathcal X/U)
\]
のperfect bimoduleのmoduliにToën--Vaquié \cite{TV} を適用する。
同値を与えるbimoduleはunit/counitのconeの消滅で定義される開部分であり、
その上にrotation intertwinerを付加する。
負Extの消滅によって、必要な部分は通常のArtin 1-stackになる。

指定された還元同値を保ってsmall extensionに沿うliftを作ることが、
「同型類が形式的に一致する」という主張より一段強い入力である。
本文ではこの点を独立補題とし、homotopy fibre内の差を幾何学的接ベクトルで消している。
局所有限表示のsmooth atlasに移してから、Stacks Project Tag 02HXと055U \cite{Stacks}を使う。
これによりcubic、bimodule、intertwinerを同じétale近傍上で得る。

\subsection{量子化族へまだ自動適用できない理由}

論文第I部の量子化族の各Artin切断でsmooth properな圏が得られることと、
開集合 \(U\) 全体上で一つのperfectな自己Hom代数と対角を得ることは別である。
後者が得られれば、可換化はétale局所実現定理へ帰着できる。
現時点では、\(D(\QGr B_U)^c\) の有限表示の代数的模型と
relative tails cohomologyの一様な制御が未完である。

étale近傍の像がZariski開なので、既約底で一点実現すればdense openの各点にモデルを得る。
ただし、単一のcubic族の元の底への降下、同値のmonodromy、開部分の外の全点での存在は、
この証明で解いたことにならない。

\subsection{四次元以外でrotationを省く意味}

\(n\ne4\) では \(\HH^2=J_3\) なので、同じArtin帰納法を
\(CC^\bullet(\Ku(X))[2]\) 自身に適用できる。
固定した変形圏上の二つのrotation liftの比は恒等に還元されるinvertible bimoduleであり、
\(\HH^1=0\) により対角と同型となる。
詳細は論文の「四次元以外ではrotationを忘れた形式復元」に置く。

full higher groupoidでは、pairの接複体の \(H^{-1}\) にscalar intertwinerの
\(\CC\) が現れる一方、rotationなしではその項がない。
従って同型類の全単射をhigher groupoidの同値へ読み替えない。
\(n=7\) で \(\HH^3\ne0\) であっても、存在するliftどうしの差を消す方法は使える。

\subsection{四次元の一般非可換変形との比較に残る点}

Huybrechts \cite[Section 9]{HuyUpdate} は
\(\HH^2(X)=\CC^{20}\oplus\CC^{20}\) と
\(\HH^2(\Ku(X))\simeq\CC^{22}\) の関係を取り上げ、
全体の非可換変形が残余K3圏の余分な方向を説明する可能性をHodge理論から示唆する。
これはその節で証明された制限定理ではない。

一方、本稿のgraded変形から作るpairでは、制限像はrotation不変な20方向へ入る。
「gradedという仮定が強いから両者は自動的に整合する」とだけ説明するのは不十分である。
論文の表示 \(q=\iota_{dF}P\) に従えば、任意の定数trivector由来のbivectorについて
\[
 f*g=fg+\frac\varepsilon2q(f,g),\qquad\varepsilon^2=0
\]
は一次の結合的なgraded変形となり、\(F\) がCasimirなので商にも降りる。
一次ではJacobi条件自体は不要である。
従って、HKR分解、残余射影によるHochschild制限、次数シフトの変形を
同一の模型で比較することが有用な次の課題である。
現時点ではこの比較を完了しておらず、既知定理との矛盾や、追加方向の消滅を断定しない。

\section{相対Beilinson模型から得られるproperな族}
\label{sec:relative-proper-model}

本節は、一般の量子化族の相対saturationに向けた部分結果である。
以下の相対Beilinson仮定のもとで、超曲面圏と残余圏を
実在する基底上のproper dg代数で表す。
この結論にsmoothness、homotopical finite presentation、
または一般の有限パラメータでの可換復元を含めない。
相対Beilinson仮定そのものと、最後に残る対角加群の有限性は別の問題である。

\subsection{仮定と基底変更の意味}

\(R\) をNoetherな可換 \(\CC\)-代数、\(N=n+2\ge5\) とする。
\(A=\bigoplus_{d\ge0}A_d\) は \(A_0=R\) を満たす次数付き代数で、
有限生成の \(A_1\) により生成され、各 \(A_d\) は有限射影 \(R\)-加群とする。
中心元 \(\Omega\in A_3\) は非零因子で、
\(B=A/(\Omega)\) の各次数成分も有限射影 \(R\)-加群とする。
したがって
\[
 0\longrightarrow A(-3)\xrightarrow{\Omega}A
   \longrightarrow B\longrightarrow0
 \tag{Q}\label{eq:relative-quotient-sequence}
\]
は任意の可換環への底変更後にも完全である。
\(\QGr\) は全次数付き右加群の圏をtorsion加群で割ったGrothendieck圏、
\(\OO_A(i)\), \(\OO_B(i)\) は構造加群のシフトの像を表す。

\begin{definition}[本節の相対Beilinson仮定]\label{def:relative-beilinson-hyp}
 必要とする各Noetherな底変更 \(R\to R'\) に対し、次を仮定する。
 \begin{enumerate}[label=(RB\arabic*)]
 \item 次数付き加群の導来局所化は
 \[
 D(\QGr A_{R'})\simeq
 D(\operatorname{Gr}A_{R'})/
 D_{\operatorname{Tors}}(\operatorname{Gr}A_{R'}),
 \]
 および \(B_{R'}\) に対する同じ表示を持つ。
 これらのdg enhancementを用い、次数シフトと導来係数拡大が
 局所化へ降りるものとする。
 \item \(G_{R'}=\bigoplus_{i=0}^{N-1}\OO_{A_{R'}}(i)\) はcompact generatorで、
 \[
 \REnd(G_{R'})\simeq T_{R'},\qquad
 e_jTe_i=A_{j-i}\quad(0\le i,j\le N-1),
 \tag{B}\label{eq:relative-beilinson-algebra}
 \]
 である。ただし \(A_d=0\) for \(d<0\) とし、右辺は次数零の三角代数である。
 積はHomの合成によるものとし、この側のBeilinson代数を \(T=\nabla A\) と書く。
 oppositeに関する行列の規約はこの等式で固定する。
 同値は自然な \(T\otimes_RR'\to\REnd(G_{R'})\) によって与えられる。
 \end{enumerate}
\end{definition}

ここで(RB2)は、体上のAS-regular性を述べる条件より強い。
各fiberにおけるBeilinson型定理を引用するだけでは、
この普遍的な底変更を含む条件を検証したことにはならない。
通常の相対構成では、有限長の線形Koszul分解、
その双対複体、torsion局所化の計算を用いて(RB2)を検証する。
本節ではこの検証を、任意の量子化族について完了したとは主張しない。

\subsection{超曲面圏のproper模型}

\begin{proposition}[相対proper模型]\label{prop:relative-proper-model}
 定義\ref{def:relative-beilinson-hyp}の条件のもとで
 \[
 E=\bigoplus_{i=0}^{N-1}\OO_B(i),\qquad
 \Lambda=\REnd_{\QGr B}(E)
 \]
 とおくと、\(E\) は \(D(\QGr B)\) のcompact generatorであり、
 \[
 D(\QGr B)^c\simeq\Perf(\Lambda)
 \]
 を得る。\(\Lambda\) は \(R\)-perfectなdg代数である。
 また、仮定の範囲の各底変更に対して、dg代数として
 \[
 \Lambda\otL_RR'\simeq
 \REnd_{\QGr B_{R'}}(E_{R'})
 \tag{BC}\label{eq:relative-proper-basechange}
 \]
 が成り立つ。
\end{proposition}

\begin{proof}
 次数付き加群のrestrictionはtorsion加群をtorsion加群へ送り、
 \(B\)-加群について \(A\)-torsionと \(B\)-torsionは一致する。
 よってrestrictionはexact、conservativeかつ直和を保つ関手
 \[
 j_*:D(\QGr B)\longrightarrow D(\QGr A)
 \]
 を誘導する。その左随伴 \(j^*\) は導来係数拡大である。
 実際、\eqref{eq:relative-quotient-sequence}によって
 \(-\otL_A B\) はtorsion cohomologyを持つ複体を同種の複体へ送り、
 導来局所化へ降りる。随伴のunitとcounitも降りる。
 また、restriction後の \(j_*j^*\) には自然なtriangle
 \[
 X(-3)\xrightarrow{\Omega}X\longrightarrow j_*j^*X
 \tag{M}\label{eq:relative-quotient-monad}
 \]
 がある。\(j_*\) の直和保存とconservativityから、
 \(j^*G=E\) はcompact generatorとなる。

 \(T=\nabla A\) は有限射影 \(R\)-加群である。
 さらに \(S=R^{\oplus N}\subset T\) とすると、
 augmentation \(T\to S\) の核 \(J\) は狭義三角行列からなる。
 \(J^{\otimes_S q}=0\) for \(q\ge N\) であり、
 normalized relative bar resolution
 \[
 \cdots\longrightarrow
 T\otimes_SJ^{\otimes_S2}\otimes_ST
 \longrightarrow T\otimes_SJ\otimes_ST
 \longrightarrow T\otimes_ST\longrightarrow T
 \]
 は有限長で、各項は有限射影 \(T\)-bimoduleである。
 したがって \(T\) は \(R\) 上smooth properである。

 \(D(\QGr A)\simeq D(T)\) のもとで次数シフト \((-3)\) を表す
 invertible \(T\)-bimoduleを \(M_{-3}\) と書く。
 これは \(T\) のsmoothnessとMorita自己同値により
 \(\Perf(T^{\mathrm{op}}\otL_RT)\) に属し、特に \(R\)-perfectである。
 \eqref{eq:relative-quotient-monad}を \(G\) に適用すると、基底上の複体として
 \[
 \Lambda\simeq
 \operatorname{cofib}\bigl(M_{-3}\xrightarrow{\Omega}T\bigr).
 \tag{C}\label{eq:relative-perfect-cone}
 \]
 よって \(\Lambda\) は \(R\)-perfectである。

 このcofiberの表示だけから積を定義してはいけない。
 \(\Lambda\) の積と単位は実際の \(\REnd_B(E)\) のものを用いる。
 同じ内容をMorita模型で述べると、\(j_*j^*\) は
 unitとmultiplicationを備えたmonadであり、
 \eqref{eq:relative-perfect-cone}のcofiberはそのmonadを表す
 \(T\)-bimoduleの代数対象である。
 この代数構造が、随伴による
 \(\REnd_B(j^*G)\simeq\RHom_A(G,j_*j^*G)\) の合成と一致する。

 最後に、(RB2)により周囲のMorita同値は底変更と両立する。
 次数シフト、中心元 \(\Omega\)、随伴のunitとmultiplicationも
 導来係数拡大から作られるので底変更と両立する。
 したがって、monadの代数対象としての底変更から
 \eqref{eq:relative-proper-basechange}が得られる。
 これは基底複体だけの同一視より強い主張である。
\end{proof}

\begin{remark}[有限性の範囲]\label{rem:relative-model-finiteness}
 命題\ref{prop:relative-proper-model}は、各Artin切断を別々に構成して
 逆極限を取る主張ではない。\(R\) 上に存在するadjunctionとmonadから、
 \(R\)-perfectな実際のdg代数を与える。
 ただし、基底上のperfect複体に移した \(A_\infty\) 積が
 有限個の高次演算だけで記述できることや、\(\Lambda\) が
 homotopically finitely presentedであることは従わない。
 特に、この命題をsmoothnessの証明と読み替えることはできない。
\end{remark}

\subsection{局所cohomologyによる確認可能な数値条件}

(RB2)を検証する際の具体的な目標は、ある可逆 \(R\)-加群 \(L\) に対する
次の構造層cohomologyである：
\[
 H^pR\Gamma(\OO_A(d))\simeq
 \begin{cases}
 A_d,&p=0,\ d\ge0,\\
 A_{-N-d}^{\vee}\otimes_RL,&p=N-1,\ d\le-N,\\
 0,&\text{otherwise}.
 \end{cases}
 \tag{H}\label{eq:relative-ambient-cohomology}
\]
この式そのものにも底変更との両立を要求する。
有限線形Koszul分解の存在はshiftの有限生成性を、
\eqref{eq:relative-ambient-cohomology}はtiltingのExt消滅を与える。
一方、compact性と全生成性を得るには、torsion局所化の導来関手が
必要な直和を保つこと、または同等のcompact generationの証明が別途必要である。

\eqref{eq:relative-ambient-cohomology}が分かれば、
\(0\le i,j\le N-1\), \(d=j-i\) に対して
\[
 H^p\RHom_B(\OO_B(i),\OO_B(j))\simeq
 \begin{cases}
 B_d,&p=0,\ d\ge0,\\
 A_{3-N-d}^{\vee}\otimes_RL,
     &p=N-2,\ 0\le3-N-d\le2,\\
 0,&\text{otherwise}
 \end{cases}
 \tag{H$_B$}\label{eq:relative-hypersurface-cohomology}
\]
が \eqref{eq:relative-quotient-monad}から従う。
特にこの窓のHomは基底上有限射影cohomologyを持ち、その振幅は
\([0,N-2]=[0,n]\) に含まれる。
この式は積の全情報を代替しない。

\begin{lemma}[fiberwise Koszul性から得られる有限分解]
\label{lem:relative-koszul-resolution}
 上の次数別有限射影性を仮定し、各幾何学的fiber \(A_s\) が
 Koszulで、augmentation加群 \(k(s)\) の射影次元が \(N\)、
 Hilbert級数が \((1-t)^{-N}\) であるとする。
 このとき有限射影 \(R\)-加群 \(W_p\) を用いる普遍的に完全な線形分解
 \[
 0\to W_N\otimes_RA(-N)\to\cdots\to
 W_1\otimes_RA(-1)\to A\to R\to0
 \tag{K}\label{eq:relative-linear-resolution}
 \]
 が存在し、\(\rank W_p=\binom Np\) である。
\end{lemma}

\begin{proof}
 \(A_+\) の次数別成分は有限射影である。
 \(W_1=A_1\) として、射影加群 \(W_1\otimes_RA(-1)\) から \(A_+\) への射を取る。
 各fiberで全射なので、各次数でNakayamaの補題により全射である。
 像の次数別成分が射影であるため、核も次数別有限射影であり、
 核の形成は任意の底変更と両立する。
 同じ操作をsyzygyに反復し、\(p\) 次syzygyを生成する次数 \(p\) の成分を
 \(W_p\) とする。fiberのKoszul性により、この成分がsyzygy全体を生成する。
 したがって各段階で同じ射影性と底変更の主張が成り立つ。
 第 \(N\) 段階の次の核はすべてのfiberで零なので、次数ごとのNakayamaにより零である。
 fiberのKoszul Betti数はHilbert級数から \(\binom Np\) となる。
 各次数の短完全列は \(R\)-加群として分裂するので、分解は普遍的に完全である。
\end{proof}

この補題は有限なKoszul分解を構成するが、
各fiberのAS-regular性だけからrelative tailsの性質をすべて得たという主張ではない。
双対Koszul複体と導来torsion局所化を含めて
\eqref{eq:relative-ambient-cohomology}および(RB2)を検証する必要が残る。

\subsection{残余圏のpropernessとrotationの開性}

さらに \(\OO_B,\ldots,\OO_B(N-4)\) が相対例外列で、
そのHomと底変更が両立すると仮定する。
例えば \eqref{eq:relative-hypersurface-cohomology}はこの条件を与える。
\[
 \cK_R=\langle\OO_B,\ldots,\OO_B(N-4)\rangle^\perp
 \subset\Perf(\Lambda)
\]
とおく。

\begin{proposition}\label{prop:relative-residual-proper}
 \(\cK_R\) のinclusionは左随伴 \(p\) を持ち、
 \(G_{\cK}=p(E)\) はcompact generatorである。
 \(\Gamma_R=\REnd(G_{\cK})\) は \(R\)-perfectで、
 \(\cK_R\simeq\Perf(\Gamma_R)\) は底変更と両立する。
 また \(\rho=L_{\OO_B}\circ(1)\) は \(\cK_R\) の自己関手を与え、
 底変更と両立する。
\end{proposition}

\begin{proof}
 propernessにより各相対例外対象 \(F\) に対し、
 \(\Perf(R)\to\Perf(\Lambda)\), \(V\mapsto V\otimes_RF\) は
 左右随伴 \(\RHom(-,F)^\vee\), \(\RHom(F,-)\) を持つ。
 その左mutationは \(F^\perp\) への左projectionを与える。
 \(a=N-3\), \(F_i=\OO_B(i)\) とおくと
 \[
 p=L_{F_0}\circ L_{F_1}\circ\cdots\circ L_{F_{a-1}}
 \]
 は \(\cK_R\) への左projectionである。
 実際、右側から順に適用し、\(j>i\) に対する
 \(\RHom(F_j,F_i)=0\) を用いると、既に得た直交性が保存される。
 各evaluation triangleはcompact対象の有限なtriangleであり、
 \(p|_{\cK_R}=\Id\) なので、\(p(E)\) は残余圏を生成する。
 そのEnd複体のpropernessおよび底変更は、
 元のproper dg模型とこれらのevaluation triangleから従う。

 \(F\in\cK_R\) とする。\(1\le j\le a-1\) なら
 \(\RHom(\OO_B(j),F(1))=\RHom(\OO_B(j-1),F)=0\) であり、
 \(\RHom(\OO_B(j),\OO_B)=0\) である。
 \(L_{\OO_B}\) のevaluation triangleから \(\rho(F)\in\cK_R\) が従う。
 同じtriangleから底変更との両立も従う。
\end{proof}

\begin{remark}[両側admissibilityは追加条件]
 proper性だけから \(\cK_R\) のinclusionが右随伴も持つとは結論しない。
 同様に、\(L_{\OO_B}\) または \(\rho\) が右随伴を持つことも、
 上の証明だけからは従わない。
 例外対象への関手の左右随伴が存在することと、
 その右直交圏の両側admissibilityは別の主張である。
\end{remark}

\begin{lemma}[随伴が与えられた場合のrotationの開性]
\label{lem:relative-rotation-open-conditional}
 上の状況で、さらに \(\rho\) がcompact対象上の \(R\)-linear右随伴
 \(\rho^R\) を持ち、そのunitとcounitを含めて底変更と両立すると仮定する。
 ある点 \(s\in\Spec R\) で \(\rho_s\) が自己同値ならば、
 \(s\) のZariski近傍へ制限することで \(\rho\) は自己同値となる。
\end{lemma}

\begin{proof}
 unitとcounitを生成子 \(G_{\cK}\) 上で評価した二つのconeはcompact対象である。
 これらからの \(\RHom(G_{\cK},-)\) は \(R\)-perfectで、
 fiber \(s\) で零になる。
 perfect複体のsupportの閉性により、両方が零となる \(s\) の開近傍がある。
 そこで生成子性を用いると二つのconeは零であり、
 unitとcounitが生成子上、従ってそのthick閉包全体で同型となる。
\end{proof}

上の右随伴の追加仮定は、例えば底変更と両立するrelative Serre dualityを
周囲の圏に別途与えれば検討できるが、本節のproper模型からは導かない。
またrotationの可逆性が分かった場合でも、そのkernelが
\(\Gamma_R^{\mathrm{op}}\otL_R\Gamma_R\) 上perfectであることは別の条件である。
対角加群のperfect性が未証明なら、恒等関手についてさえこの条件は未証明である。

\subsection{smoothnessに十分な有限の証明書と残る問題}

\begin{lemma}[対角の有限証明書]\label{lem:relative-diagonal-certificate}
 \(\Gamma\) を \(R\)-perfectなdg代数とし、\(\Gamma^e=\Gamma^{\mathrm{op}}\otL_R\Gamma\)
 と書く。ある \(P\in\Perf(\Gamma^e)\) と射 \(P\to\Gamma\) が存在し、
 点 \(s\) での導来fiberがquasi-isomorphismならば、\(s\) の開近傍で \(\Gamma\) はsmoothとなる。
\end{lemma}

\begin{proof}
 \(\Gamma^e\) は \(R\)-perfectなので、有限cellと直和因子で作られる \(P\) も
 \(R\)-perfectである。\(\Cone(P\to\Gamma)\) も \(R\)-perfectであり、
 \(s\) でのfiberが零なので、ある開近傍で零である。
 従ってその近傍では対角 \(\Gamma\simeq P\) が \(\Gamma^e\) 上perfectである。
\end{proof}

\begin{remark}[現時点で未解決の点]\label{rem:relative-saturation-gap}
 fiber \(\Gamma_s\) のsmoothnessは、そのfiber上での有限cell対角分解を与える。
 しかし、その有限cell、微分、retract、および対角への射を
 \(R\) の近傍へ同時に持ち上げた証明書はまだ構成していない。
 各Artin切断でのsmoothnessだけではこの構成を代替できない。
 Toën \cite[Lemma 3.14]{ToenFinitude} が述べるのは、properなdg代数の
 「smoothかつcohomological dimension \(\le d\)」という部分前層の
 quasi-compactnessであり、一般の開性ではない。
 この補題をsmoothnessの開性として使わない。

 また、Toën--Vaquié \cite[Theorem 3.6]{TV} の有限型dg圏のobject moduli定理を、
 proper性だけが分かっている \(\Gamma\) に適用することはできない。
 特にpseudo-perfect加群とcompactなperfect加群の同一視にはsmoothnessが必要である。
 既知の飽和圏内で生成子を動かして比較する方法にも、
 生成性の開性および比較moduliの有限表示性の証明が別途必要となる。
 これらのmoduliによる迂回を完成した代数化定理として扱わない。
\end{remark}


\subsection{有限例外列による相対smoothnessの還元}
現行論文の「相対saturationへの二つの還元」節の命題により、
相対properなdg圏$\mathcal C_R=\Perf(\Lambda_R)$が
有限個の相対例外対象と残余圏$\mathcal K_R$に分解するとき、
$\mathcal C_R$と$\mathcal K_R$の相対smoothnessは同値となる。
gluing bimoduleは$R$-perfectで、残余圏がsmoothなら
対角のperfect性により実際にperfect bimoduleとなる。
従って求める証明書は$\Lambda_R$の対角について一つあればよい。
これは(RB1)--(RB2)を仮定する還元であって、相対smoothnessそのものの
証明ではない。閉ファイバーのBeilinson対角分解を全族へ
持ち上げる問題が残る。詳細な証明は現行論文
の同節に置く。

graded matrix factorizationの内部シフトは
中心元の次数が三ならば常に$(1)^3\simeq[2]$を満たす。
これは相対Orlov同値が与えられればrotationの自己同値性と
相対三次周期性を保証する。ただし相対Orlov同値自体は
依然として未構成であり、形式的周期性からその存在は導けない。

\section{一ブロックの有限パラメータ比較}

\subsection{中心元の正規化}

\[
 T=\frac16(xyz+xzy+yxz+yzx+zxy+zyx),\qquad G_h=D-2h^2T,
 \quad D=x^3+y^3+z^3.
\]
\(T=xyz+(h/2)(-x^3+y^3-z^3)\) なので
\[
 G_h=(1+h^3)(x^3+z^3)+(1-h^3)y^3-2h^2xyz.
\]
Waltonの標準中心元 \cite[Equation (2.16)]{Walton} はこの規約で \(-hG_h\) になる。
この符号と \(h^2\) を混同するとHesse曲線の係数を誤る。
点曲線は \(D+h^2xyz=0\) であり、\(D-2h^2xyz=0\) ではない。

非零性は、三つのcubeに1、三つの巡回語に \(h/2\)、逆巡回語に \(-h/2\)、
二文字だけが等しい語に零を取る汎関数 \(\ell\) で確認する。
これは全18個の次数三の関係式を殺し、\(\ell(G_h)=3\) となる。
特殊値で中心の次元が増えても、この非零性の証明は変わらない。

\subsection{楕円曲線のtranslationとrotation}

右加群のTHCR規約では、tailsの次数シフトは
\(H(F)=\sigma_*(F\otimes L)\) である。
windowを一つ動かすと
\[
 \rho_{\rm nc}=T_{\OO_E}H=t_{-p}^*\rho_0,
 \qquad \rho_0=T_{\OO_E}(-\otimes L),\quad\sigma=t_p.
\]
\(A_q=(-\otimes\OO_E(q-o))\) とすると、Fourier--Mukai核のsee-sawから
\[
 A_q\rho_0A_q^{-1}=t_q^*\rho_0.
\]
従って \(A_{-p}\) で \(\rho_{\rm nc}\) を \(\rho_0\) に共役できる。
periodicityの比較の差はscalarであり、三乗根を選んでintertwinerを修正する。
この修正まで行ってから有限次数差の群作用を移す。

\subsection{任意の中心三次式と次数の降下}

中心三次式 \(P(u)\) が原点だけに臨界点を持てば、Jacobian ideal \(J(P)\) は
\((u)\)-primaryである。factorization微分を \(u_i\) で微分すると、
\(\partial_iP\) の零ホモトピーを得る。
そのため、適当な \(u_i^d\) のKoszul coneへ支持を縮め、有限filtrationによって
\(S/(g)\)-加群の制限へ生成性を帰着できる。

Euler公式 \(P=\sum u_iq_i\)、\(q_i=\frac13\partial_iP\) に対する
Koszul factorizationを用いれば、制限加群の安定化を外部積として書ける。
内部次数は最初に
\[
 \Gamma=(\ZZ e_0\oplus\ZZ e_1)/(3e_0-3e_1)
\]
を用い、全中心変数を \(e_1\) 次数とする。
\(\ker(\Gamma\to\ZZ)\simeq\ZZ/3\) の双対群が、全中心変数の同時三乗根倍に対応する。
rotationとperiodicityを保つブロック同値を、双対同変化で通常次数へ降ろす。

Fermatの場合は各 \(u_i\) に独立な次数を与えた
\(\Gamma_m=\ZZ^{m+1}/\langle3e_i-3e_0\rangle\) も使える。
この細分次数は多ブロックの完全忠実比較で用いる。
例外値のnonpropernessは、二ブロック次数と非零なKoszul factorizationだけでも検出できる。
一般の \(P\) に対して、個々の \(u_i\) の三乗根倍が \(P\) を保つとは仮定しない。

\subsection{初期cubicが固定されるわけではない}

明示モデルは
\[
 F_h=x^3+y^3+z^3+\sum u_i^3+h^2xyz.
\]
FermatのJacobian ringは全変数の二乗を零にするので \([xyz]\ne0\) である。
従って一次の \(h\)-項が消えても、二次の幾何学的変形は非零である。
\(h\)-形式族として初期Fermat cubicに定数化できる、という結論にはならない。
ただしこれは本稿で固定した中心元の族についての記述である。

\section{一つの無限位数因子を許す多ブロックと生成性の障害}

\subsection{有限中心性で解決する段階}

\(C_j=S_{h_j}/(G_{h_j})\) と置く。
Euler導分と中心変数のJacobian零ホモトピーは、全MF圏の生成性を
\(C=\bigotimes_j C_j\) 上の有限生成加群のrestrictionへ帰着する。
その後に必要なのは
\[
 D^b(\operatorname{gr}^{\Gamma}C)
 =\thick\{\operatorname{coarsen}(M_1\boxtimes\cdots\boxtimes M_r)(\gamma)\}
\]
である。右辺のconeを作る射は、粗化後の圏の全ての次数零の射を使う。
各対象や各射の \(\ZZ^r\)-次数への持ち上げを要求してはいけない。

論文の「一方有限中心の外部積生成性と多ブロックの可換モデル」では、
任意の左右noetherianな \(\Lambda\)-次数付き係数代数 \(A\) に対し、
第二因子 \(D\) だけが有限型可換斉次中心上有限なら
\[
 D^b(\operatorname{gr}^{\Lambda}(A\otimes D))
 =\thick\{M\boxtimes\operatorname{coarsen}_\beta N\}
\]
を証明する。\(A\) の有限中心性や有限global dimensionは不要である。
詳細な証明は論文の「非可換係数を持つ多項式拡大」の補題、
「多項式中心上の分離的な稠密開集合」の補題、
「一方だけが有限中心である外部積生成」の定理にまとめた。
検証に必要な要点は以下である。

\paragraph{非可換係数のproper pushforward.}
\(A\otimes k[t_1,\ldots,t_d]\)-加群を、\(\PP^d\) 上の
\(A\otimes\OO\)-coherent加群へ斉次化する。
これは \(\OO\)-coherentである必要はない。
\v Cechの長さと \(A\otimes\OO(a)\) からの全射を使う降下帰納法により、
導来大域切断のcohomologyが有限生成 \(A\)-加群であることを示す。
有限global dimensionに置き換えないことが重要である。
Beilinson分解を \(t_0\ne0\) に制限して多項式拡大の生成性を得る。

\paragraph{分離性の役割.}
中心のNoether正規化 \(P\subset Z\) を用い、斉次元を反転すると、
冪零イデアルを除いた商 \(Q\) は \(P_f\) 上有限射影かつ分離的になる。
\(A\otimes Q\)-加群 \(T\) は \(\operatorname{Res}T\otimes_{P_f}Q\) の直和因子である。
分離冪等元の元の整数次数0成分を取り、sectionが粗化次数を保つことを確認する。
\(A\otimes Q\) が有限global dimensionを持つ、という誤った帰結は使わない。

\paragraph{閉集合の回収.}
\(f\)-torsionの有限filtrationは \(D/(f)\) の加群へ帰着する。
その有限中心の次元が下がるので帰納できる。局所化の核を含めることで軸torsionも失わない。
元の \(\ZZ^r\)-次数へ持ち上がらない対象を全て含む。

Elagin--Lunts--Schn\"urer \cite[Theorem 4.18]{ELS} は双方をcoherentな構造代数とする
ungradedな生成定理を与える。ここで任意のnoetherian係数 \(A\) を許すには、
上記の係数付きBeilinson議論を別に証明する必要がある。
一方有限中心という定式化の先行研究上の新規性は未確認である。
He--Ma--Ye \cite{HMY} は中心二次元の和をClifford変形で扱い、
本稿の三次元の和にはそのまま適用できない。

\subsection{有限位数条件の正確な使用箇所}

\(\sigma_j\) の位数を \(q_j\) とすると、THCRの \(q_j\)-Veroneseは
\(\sigma_j\)-不変な正次数divisorに対応する可換section ringになる。
その有限群不変部分は元のTHCRの中心に含まれ、THCRはこの中心部分代数上有限である。
論文にはskew Laurent表示を用いた直接証明を置く。

有限中心性から \(\bigotimes_jC_j\) はnoetherianである。
正次数の中心正則元 \(G_j\) を逆順に戻すと、周囲のtensor productもnoetherianとなる。
Koszul双加群分解のtensor productはAS-regular性を与える。
この方法でもMF圏をqgrのOrlov残余成分と比較するための周囲の仮定を確認できる。
さらに本文の一般の楕円THCRの積の補題は、有限位数の条件なしで同じnoetherian性を与える。

\paragraph{多重amplenessの先行研究との関係.}
Keeler \cite[Theorem 2.7, Corollary 2.8]{KeelerMulti} は、任意有限個の可換な可逆bimoduleについて、
数値的Picard群への作用のquasi-unipotenceと十分正の多重冪のamplenessから、
左右のNC-amplenessを判定する。本研究の積のtranslationは数値的に自明に作用し、
正次数線束の正の外部積はampleなので、この判定にも適合する。
一方、同論文のHypothesis 3.1を仮定するTheorem 3.3と、THCRの積に関するCorollary 3.5は
二因子の主張である。任意個の因子を扱うSection 4のHypothesis 4.1は、
代数の右noetherian性を既に仮定している。
従って本文では固定面と有限幅の帯の帰納法を明示し、任意個の積のnoetherian性を証明した。
この特殊な多因子への拡張の新規性は主張しない。
参照番号はarXiv一次TeXソースで確認したもので、刊行版との番号一致は未照合である。

可換モデルは
\[
 V\!\left(\sum_j(x_j^3+y_j^3+z_j^3+h_j^2x_jy_jz_j)+P(u)\right).
\]
\(P\) は原点だけに臨界点を持つ任意の中心三次式でよく、中心変数なしも許す。
主定理は \(j\ge2\) のみを有限位数に制限し、\(C_1\) は任意位数のnoetherian係数として残す。
最初の \(C_1\)-加群は \(\Gamma/\ZZ e_1\) の剰余類ごとに分解でき、
以後は一方有限中心定理を反復する。
この開始時の分解を、積の一般の加群が細分次数を持つという主張に拡大してはいけない。
全てのパラメータを同時に任意にする定理ではなく、二つ以上の無限位数因子は未解決である。

\subsection{非零パラメータの厳密な証明書}

Hesse曲線の原点を \(o=[1:-1:0]\)、平行移動点を \(p=[1:-1:-h]\) と取る。
論文の位数5の補題では
\[
 2p=[h^3+1:h^3-1:-2h],\qquad
 T_{2p}(p)=h^3(5-h^6)
\]
を使う。\(T_{2p}\) は、表示された座標で曲線のgradientが定める接線式である。
\(h^6=5\) なら、接線の交点の和が零になる群法から \(5p=0\) となり、
\(p\ne o\) と5が素数であることから位数はちょうど5である。
これは近似計算による位数の推定ではない。

\texttt{computations/verify\_algebra.py} は、点の曲線への帰属、倍点の接線条件、
接線とcubicの交点重複度、および \(T_{2p}(p)\) の恒等式を有理係数で確認する。
群法から位数を結論する論証は論文に置く。
たとえば正の実数 \(h=5^{1/6}\) を各ブロックに選べば、モデルの混合項の係数は \(5^{1/3}\) となる。
\((r,m)=(2,2)\) と \((3,0)\) はそれぞれsixfoldとsevenfoldを与える。

\paragraph{無限位数の証明書.}
\(h=2\) では \(E_2:x^3+y^3+z^3+4xyz=0\) と \(p_2=[1:-1:-2]\) は有理数体上定義される。
5と11で良い還元を持ち、点数はそれぞれ9と12である。
prime-to-residue-characteristic torsionの単射性から、もし \(p_2\) がtorsionなら位数は3を割る。
しかし \(2p_2=[9:7:-4]\ne[1:-1:2]=-p_2\) なので不可能である。
計算コードの点数と射影座標の確認が有限の証明書になっている。
\(h_1=2\)、他の \(h_j^6=5\) は、有限位数のみの族には含まれないsixfold・sevenfoldの例である。

\subsection{二つ以上の無限位数で再使用できること・できないこと}

無限位数でも、ブロックごとの楕円圏の同値、外部積の完全忠実性、
Euler零ホモトピー、周囲のnoetherian AS-regular性は保たれる。
一つの無限位数因子は係数代数として残せるが、二つ目を追加する際には
一方有限中心定理の仮定を失う。
従って最小の未解決例は \(r=2,m=0\) で両方の平行移動が無限位数の場合である。
まずその \(C_1\otimes C_2\) の粗化次数付き外部積生成に集中する。
証明できれば、追加の有限中心ブロックと任意の孤立中心三次式は現行の議論で処理できる。
任意個の無限位数へ進むには、任意noetherian係数を許すTHCR版の生成補題が十分である。

注意すべき対象の例は、\(3e_1=3e_2\) の粗化次数を持つ
\(\CC[a,b]/(a^3-b^3)\) である。
そのannihilatorは \(\ZZ^2\)-斉次でないので、同じ加群の細分次数へのliftは存在しない。
しかし粗化後には \(a^3-b^3\) による二つのfree shiftの間のconeで生成される。
従ってこの例は全生成性の反例ではなく、誤ったlift仮定を防ぐための例である。
multigraded tailsだけを使う方法は、軸上のtorsionを別に回収する必要がある。
有限位数の稠密性だけから、二つ以上の無限位数因子への拡張を結論してはならない。

\section{PIの中心幾何から得られる有限生成子}
\label{sec:notes-pi-residue-generator}

有限位数のSklyaninブロックでは、本文の外部積生成定理とは別に、
中心の零sliceのAzumaya幾何から全MF圏のpropernessと有限生成子を直接得られる。
この証明は \(\gcd(|\sigma_i|,3)=1\) を仮定しない。

\paragraph{設定.}
\(r\ge1\)、\(m\ge0\)、\(h_i^6\ne-27\) とし、
\[
 \begin{aligned}
 A&=\bigotimes_{i=1}^rS_{h_i}\otimes_\CC T,
 &T&=\CC[u_1,\ldots,u_m],\\
 W&=\sum_{i=1}^r g_i+P,
 &R&=A/(W)
 \end{aligned}
\]
と置く。ここで \(g_i=G_{h_i}\)、\(P\in T_3\) であり、
\(m>0\) の場合は \(T/(\partial_1P,\ldots,\partial_mP)\) が有限次元、
\(m=0\) の場合は \(T=\CC,P=0\) とする。
全ての \(\sigma_i\) は有限位数 \(q_i\) を持つものとする。
\(h_i=0\) は多項式環のfactorで、\(q_i=1\) と解釈する。
周囲の代数 \(A\) はnoetherian AS-regularであり、\(W\) は中心正則元である
（本文のテンソル積の正則性の命題）。
従ってMFと特異圏の比較 \cite[Theorem 5.8]{CCKM} を用いられる。
以下ではこれらの圏を冪等完備化し、\(k=R/R_{>0}\) は特異圏での像、
またはその有限自由stabilizationも表す。

\(s=r\) (\(m=0\))、\(s=r+1\) (\(m>0\)) と置き、
\[
 \Gamma_s=\ZZ^s/\langle3e_i-3e_1:2\le i\le s\rangle,
 \qquad w=3e_i
\]
とする。各 \(S_{h_i}\) の生成元の次数を \(e_i\) とし、
中心変数があれば全ての \(u_a\) の次数を \(e_s\) とする。
通常次数は \(\Gamma_s\to\ZZ,\ e_i\mapsto1\) による粗化であり、このとき \(w=3\) である。
この写像は有限の核を持ち、正の通常次数を与えるため、
\(\Gamma_s\)-斉次自由分解と斉次射を用いる同じMF--特異圏比較が適用できる。

\begin{proposition}[PIの場合のpropernessと残余体生成]
\label{prop:notes-pi-residue-generator}
上の仮定の下で \(\MF^{\Gamma_s}(A,W)\) と \(\MF^\ZZ(A,W)\) はproperである。
特異圏において
\[
 \begin{aligned}
 \Dsg^{\Gamma_s}(R)&=\thick(K_{\Gamma_s}),\\
 K_{\Gamma_s}&=
 \bigoplus_{(a_1,\ldots,a_s)\in\{0,1,2\}^s}
 k(a_1e_1+\cdots+a_se_s)
 \end{aligned}
 \tag{N1}\label{eq:notes-pi-gamma-generator}
\]
および
\[
 \Dsg^\ZZ(R)=\thick(K_{\ZZ}),\qquad
 K_{\ZZ}=k\oplus k(1)\oplus k(2)
 \tag{N2}\label{eq:notes-pi-z-generator}
\]
が成立する。
\end{proposition}

\begin{proof}
\paragraph{一次資料から用いる中心の事実.}
\(Z_i=Z(S_{h_i})\)、\(Y_i=\Spec Z_i\) とし、\(o_i\) を斉次頂点とする。
\(h_i\ne0\) のとき、\(S_{h_i}\) が有限生成可換中心 \(Z_i\) 上有限であることは
\cite[Theorem 3.7(c), p.~38]{SmithTate} による。
さらに \cite[Theorem 1.3(1),(2a)]{WWY3} は
\[
 V(g_i)\setminus\{o_i\}\subset Y_i^{\mathrm{sm}}
       =\operatorname{Az}(S_{h_i}),\qquad
 V(g_i)\setminus\{o_i\}\ \text{は滑らか}
 \tag{N3}\label{eq:notes-pi-zero-slice}
\]
を与える。原論文の三次中心元 \(g\) と \(g_i\) は非零定数倍なので、
零sliceについて同じ主張になる。
従って \(y_i\in V(g_i)\setminus\{o_i\}\) において
\[
 \widehat{(Z_i)_{y_i}}\simeq\CC[[g_i,a,b]],\qquad
 S_{h_i}\otimes_{Z_i}\widehat{(Z_i)_{y_i}}
       \simeq M_{q_i}(\CC[[g_i,a,b]]).
 \tag{N4}\label{eq:notes-pi-matrix-completion}
\]
実際、滑らかな三次元中心と滑らかな零sliceから \(dg_i\ne0\) が従い、
\(g_i\) を正則パラメータに選べる。
完全局所環上のAzumaya代数は、剰余体 \(\CC\) 上でsplitし、
行列冪等元の持ち上げによって行列環になる。
\(h_i=0\) でも、行列sizeを一とし、Fermat三次式の勾配を用いれば
\eqref{eq:notes-pi-zero-slice}--\eqref{eq:notes-pi-matrix-completion} が成立する。

\paragraph{行列Homの支持.}
\(Z=\bigotimes_i Z_i\otimes T\) は有限型可換代数であり、\(A\) は \(Z\) 上有限である。
有限自由factorization \(Q,Q'\) の次数を忘れた二周期行列Hom複体を
\[
 \overline C=\Hom_A^{\ZZ/2}(Q,Q'),\qquad
 \partial(f)=d_{Q'}f-(-1)^{|f|}fd_Q
\]
と書く。その二つの項は有限生成 \(Z\)-加群なので、
\(\overline H^\epsilon=H^{\bar\epsilon}(\overline C)\) も有限生成である。
ブロックEuler導分と中心偏微分により
\[
 d_Q\delta_i(d_Q)+\delta_i(d_Q)d_Q=3g_i\Id_Q,\qquad
 d_Q\partial_{u_a}(d_Q)+\partial_{u_a}(d_Q)d_Q
       =(\partial_aP)\Id_Q.
\]
これらをHomへ合成することで、\(g_i\) と \(\partial_aP\) は
\(\overline H^\epsilon\) を殺す。
Jacobian idealのradicalは \((u_1,\ldots,u_m)\) なので
\[
 \operatorname{Supp}_Z\overline H^\epsilon
 \subset V(g_1,\ldots,g_r,u_1,\ldots,u_m).
 \tag{N5}\label{eq:notes-pi-support}
\]

右辺の閉点 \(y=(y_1,\ldots,y_r,0)\) で、ある \(y_i\ne o_i\) とする。
他のSklyanin factorのテンソル積をその中心点で完備化した代数を \(B\) とすると、
\eqref{eq:notes-pi-matrix-completion} と中心上の完備テンソル積により
\[
 \widehat A_y\simeq
 M_{q_i}\!\left(B[[g_i,a,b,u_1,\ldots,u_m]]\right),
 \qquad W=g_i+\sum_{j\ne i}g_j+P(u).
 \tag{N6}\label{eq:notes-pi-local-matrix}
\]
ここでは次数を忘れてよい。
右辺の形式変数 \(g_i\) に関する成分ごとの導分を \(\partial_{g_i}\) とすれば、
\(\partial_{g_i}W=1\) である。
任意の完備化したfactorizationの微分 \(d\) に対して
\[
 h=\partial_{g_i}(d),\qquad dh+hd=\Id
 \tag{N7}\label{eq:notes-pi-contraction}
\]
となる。従って局所的なungraded MFは実際にcontractibleである。
行列Homにも \(\mathcal H(f)=h_{Q'}f\) による
\(\partial\mathcal H+\mathcal H\partial=\Id\) がある。
中心完備化はflatなのでHom cohomologyと可換であり、
局所環上ではfaithfully flatなので
\((\overline H^\epsilon)_y=0\) が従う。

こうして \(\overline H^\epsilon\) の支持は唯一の頂点
\(o=(o_1,\ldots,o_r,0)\) に含まれる。
\(Z\) は有限型可換 \(\CC\)-代数であるから、
\(\overline H^0\oplus\overline H^1\) は有限次元 \(\CC\)-ベクトル空間である。
この議論には、完備化した他のfactorのglobal dimensionを求める必要はない。

\paragraph{周期性とproperness.}
二周期Homを、偶数部分の微分が内部次数零、奇数部分の微分が内部次数 \(w\)
となるように次数付ける。
具体的に \(Q=(Q_1\to Q_0\to Q_1(w))\) の規約では
\[
 \begin{aligned}
 \overline C^0&=\underline{\Hom}_A(Q_0,Q'_0)
       \oplus\underline{\Hom}_A(Q_1,Q'_1),\\
 \overline C^1&=\underline{\Hom}_A(Q_0,Q'_1(w))
       \oplus\underline{\Hom}_A(Q_1,Q'_0).
 \end{aligned}
\]
とする。ここで \(\underline{\Hom}\) は全ての内部次数のHomの直和である。
\(\overline H^\epsilon\) はこの内部次数を継承し、\(\Lambda=\Gamma_s,\ZZ\) のいずれでも
\[
 \begin{aligned}
 H^{2n}\Hom_{\MF^\Lambda}(Q,Q')
      &=(\overline H^0)_{nw},\\
 H^{2n+1}\Hom_{\MF^\Lambda}(Q,Q')
      &=(\overline H^1)_{nw}.
 \end{aligned}
 \tag{N8}\label{eq:notes-pi-hom-periodicity}
\]
\(w\) は無限位数であり、有限次元の \(\overline H^\epsilon\) は有限個の内部次数しか持たない。
従ってMFのHomの全cohomologyの直和は有限次元である。
有限個のconeと直和因子にもこの結論は保たれるので、冪等完備化したMF圏はproperである。

\paragraph{残余体による生成.}
任意の有限自由factorization \(Q\) を取る。
中心の正次数斉次生成元を \(z_1,\ldots,z_t\) とすると、
全内部次数を含むstable endomorphism加群が頂点に支持されることから、十分大きい \(d\) に対して
\(z_j^d\Id_Q\) は全て零ホモトピックである。
このclassの内部次数は \(d\deg z_j\) であり、
Homの対応する斉次成分を取ることで、
\[
 z_j^d:Q(-d\deg z_j)\longrightarrow Q
\]
のnull-homotopyを次数を保った形に選べる。
中心Koszul複体とのtensor積はこれらのhomotopyを保つので、
これらの乗法のconeを順次取った対象には \(Q\) が直和因子として入る。

\(\Dsg^\Lambda(R)\) で \(Q\) を表す有限生成moduleを \(M\) とする。
今のconeの反復は、boundedな中心Koszul複体
\(K_R(z_1^d,\ldots,z_t^d)\otimes_R M\) の特異圏での像である。
そのcohomologyは \((z_1^d,\ldots,z_t^d)\) に殺される。
\(R\) は \(Z\) 上有限で、\(Z/(z_1^d,\ldots,z_t^d)\) は有限次元なので、
これらのcohomologyも有限次元である。
有限次元の次数付き \(R\)-moduleは正の通常次数によるfiltrationを持ち、
各商は有限個の \(k(\lambda)\) の直和である。
従って全ての \(Q\) は \(\{k(\lambda)\}_{\lambda\in\Lambda}\) で生成される。
冪等完備化についても同じ結論になる。

最後に \((w)=[2]\) であり、
\[
 \Gamma_s/\ZZ w\simeq(\ZZ/3)^s,\qquad \ZZ/3\ZZ\simeq\ZZ/3
\]
なので、それぞれの剰余類の代表だけで十分である。
これが \eqref{eq:notes-pi-gamma-generator} と \eqref{eq:notes-pi-z-generator} を与える。
\end{proof}

\paragraph{次の計算への使い方.}
細分次数では \(3^s\) 個、通常次数では \(k,k(1),k(2)\) の三つのstabilizationを用意すればよい。
その直和 \(K_\Lambda\) に対して
\[
 B_\Lambda=\REnd_{\Dsg^\Lambda(R)}(K_\Lambda),\qquad
 \MF^\Lambda(A,W)\simeq_{\mathrm{Morita}}\Perf(B_\Lambda)
\]
となる。自己Homはstabilization後のMFのdg模型で取り、\(H^*(B_\Lambda)\) は有限次元である。
本文の有限位数多ブロック比較と組み合わせると、この模型はsmooth properでもある。
この生成子の記述は、外部積による全生成性の証明とは独立に、
全MF圏を具体的な有限個の対象の自己Homから調べる入口を与える。
既知の中心幾何からの上の帰結について、先行研究上の新規性は未確認である。

\section{悪いパラメータと条件の必要性}

\(h^6=-27\) では点曲線が三角形になる。
\(S_h\) はAS-regular domainのままであり、\(G_h\) も非零で正則である。
商が三角形のTHCRになることは、平坦なHesse曲線の族に沿って \(G_h\) の零像を特殊化し、
三角形上のpencil trickとHilbert係数の一致で証明した。
中心の次元が一定という仮定は使わない。

そのgraded singularity categoryは \(D^b(\operatorname{coh}E_h)\) となり、
節点のskyscraperの自己Extが任意に高い次数まで残る。
ここでいうproperとは、各Hom複体が基礎体上perfect、すなわち全次数のHomの総和が
有限次元であることである。各固定次数のExtが有限というだけでは足りない。

中心変数を加えても、この無限自己Extは細分次数の外部積の中に残る。
通常次数の圏がproperなら有限群同変圏もproperとなるので、通常次数側もnonproperである。
さらに忘却後の一つの対象に任意に高い正次数の自己Homが残る。
Orlovのwindowを通して \(D^b(\qgr B)\) へ送ると、\(\qgr B\) のhomological dimensionが
無限であることが従う。有限homological dimensionなら、固定された有界複体の高次Homは消えるためである。
これはExt有限性を別途仮定しない議論である。
従って、単に比較定理の証明が失敗するのではなく、非可換射影スキーム自体がhomologicalに特異になる。

なお、斉次環 \(B\) 自身のglobal dimensionと混同しない。
そのHilbert級数の逆数は \((1-t)^{m+3}/(1-t^3)=(1-t)^{m+2}/(1+t+t^2)\) であり、
多項式ではない。従ってaugmentation加群は有限な次数付き自由分解を持たず、
\(B\) のglobal dimensionは良いパラメータでも無限である。
この判定は \(\qgr B\) のhomological dimensionとは異なる。

この例は、論文第I部の量子化に用いる開集合をさらに縮めて例外値を除くこととは両立する。
また、最初からsmooth properな非可換cubicに対象を制限した場合の反例ではない。
ただし、smooth proper性だけでも一般のrotation付き存在は従わない。
Zhang twistによる別の滑らかな反例は、周囲がCalabi--Yauの場合も含め、
論文の「可換cubic模型の存在を妨げるrotationの不変量」節に置く。
可換点を含む特定の量子化族についての全ファイバー存在は未解決である。

\section{次の問題と完了条件}

\paragraph{一般量子化の相対saturation.}
論文第I部の量子化族の残余圏とrotationを、零の近傍上のsmooth properな代数的族として作る。
これが得られればétale局所存在を適用できる。各Artin切断の存在から有限表示性を補わない。
第\ref{sec:relative-proper-model}節の相対Beilinson仮定をその量子化族で検証し、
そのproper模型について補題\ref{lem:relative-diagonal-certificate}の対角の有限証明書を構成する。

\paragraph{複数の非可換ブロック.}
高々一つの無限位数因子の場合は全生成性まで解決した。次は二つの無限位数因子の条件(G)を示す。
新たに追加するTHCRについて、一方有限中心の帰納法の代わりとなる係数付き幾何的生成を構成する。
軸torsion、粗化次数でのみ存在する対象、通常次数への降下を保持する。
独立ブロック族の周囲のnoetherian性は論文に証明があるため、核心は全生成性にある。
第\ref{sec:notes-model-boundary}節の別方法として、properness、標準Serre公式、
および \(\HH^0\) の連結性を検証することでも、既存の完全忠実比較を同値へ強められる。

\paragraph{非可換変数と中心変数の混合項.}
\(P(u)\) 内の混合項は扱えるが、非可換ブロックの関係式が \(u\) に依存する場合は別である。
\(\partial_{u_i}\) が係数代数の導分にならないため、中心変数の偏微分を用いる証明はそのまま使えない。
相対THCR、相対window、translationの相対的消去を検討する。

\paragraph{局所実現の大域化.}
étale局所モデルの元の底への降下、monodromy、実現開部分の境界を調べる。
既約底でdense openの存在が得られても、全ファイバー存在にはならない。

\section{厳密計算の役割と再現}

\texttt{computations/verify\_algebra.py} は有理係数の記号計算と有限体上の厳密な全点列挙を行う。
次を確認する。
\begin{enumerate}
 \item 汎関数 \(\ell\) に対する18個の関係式の消滅と \(\ell(G_h)=3\)。
 \item \(T\) のnormal expressionと標準中心元の \(g=-hG_h\) の恒等式。
 \item 関係行列の行列式 \(-h(D+h^2xyz)\)。
 \item Hesse曲線の特異条件 \(h^6+27=0\) の必要条件と、各例外値での特異点の構成。
 \item \(\dim J_3=\binom{n+2}{3}\)、特にsixfoldの56、sevenfoldの84。
 \item 独立な10係数を持つ一般三変数cubicに対し、外積代数上で
 \(d_K^2=s_K^2=0\)、\(d_Ks_K+s_Kd_K=P\operatorname{id}\) を有理係数で検算する。
 \item Hesse曲線の平行移動点の倍点公式と接線恒等式
 \(T_{2p}(p)=h^3(5-h^6)\)。\(h^6=5\) の位数5の論証を支える厳密な証明書である。
\end{enumerate}
\(h=2\) については \(\mathbb F_5,\mathbb F_{11}\) の全射影点を重複なく列挙し、
点数9,12、勾配の非消滅、\(p_2,2p_2,-p_2\) の射影座標の条件も確認する。
これらは本文の代数式の誤記を検出する補助であり、圏の同値、生成性、
変形理論の証明を置き換えるものではない。
計算結果は \texttt{computations/verification.json} に保存する。

\texttt{computations/verify\_rotation\_obstructions.py} は、Zhang twistのscalar関係式、
三次中心元の可換性、PBWの重なり、Nakayamaの指数、およびresidueの指標を厳密に確認する。
四次元では \(J_3\) の指標の重複度 \((8,6,6)\)、rotationの \(\HH_0\) 不変次元8と20、
不変Euler標数10と22を確認し、結果を
\texttt{computations/rotation\_obstructions.json} に保存する。
これは有限の代数・指標計算の証明書であり、Hochschildへの幾何学的同定、
同値の不存在、および連結族での局所一定性の証明は論文に置く。

組版にはLuaLaTeX、\texttt{ltjsarticle}、標準の数式・リンク関連パッケージを使う。

\begingroup
\small
\begin{thebibliography}{99}
\bibitem{PHartshorne}
R.~Hartshorne,
\emph{Algebraic Geometry}, Graduate Texts in Mathematics 52,
Springer, 1977.
\href{https://doi.org/10.1007/978-1-4757-3849-0}{doi:10.1007/978-1-4757-3849-0}.

\bibitem{PMayanskiy}
E.~Mayanskiy,
Poisson cohomology of two Fano threefolds,
\emph{Journal of Algebra} \textbf{424} (2015), 21--45.
\href{https://doi.org/10.1016/j.jalgebra.2014.08.049}{doi:10.1016/j.jalgebra.2014.08.049};
\href{https://arxiv.org/abs/1303.5916}{arXiv:1303.5916}.
本文で用いる Grassmann 記述は Theorem~1。

\bibitem{PEtingofLie}
P.~Etingof,
\emph{Lie Groups and Lie Algebras}, MIT lecture notes, 2024 course version.
\href{https://live.ocw.mit.edu/courses/18-755-lie-groups-and-lie-algebras-ii-spring-2024/mit18_755_s24_lec_full.pdf}{MIT OpenCourseWare 公開原稿}.
本文では Proposition~16.9、\S19.4、Theorem~25.17、\S32.2 を参照する。

\bibitem{PMoriSmith}
I.~Mori and S.~P.~Smith,
The classification of 3-Calabi--Yau algebras with 3 generators and 3 quadratic relations,
\href{https://arxiv.org/abs/1502.07403v2}{arXiv:1502.07403v2} (2016).
Theorems~1.3, 1.4, 4.2 and Corollaries~1.6, 2.7, 2.9.

\bibitem{PBergerTaillefer}
R.~Berger and R.~Taillefer,
Poincar\'e--Birkhoff--Witt deformations of Calabi--Yau algebras,
\emph{Journal of Noncommutative Geometry} \textbf{1} (2007), 241--270.
\href{https://arxiv.org/abs/math/0610112}{arXiv:math/0610112}.
本文で用いる potential による PBW 変形は Theorem~1.1(i)。

\bibitem{PKontsevich}
M.~Kontsevich,
Deformation quantization of Poisson manifolds, I,
\href{https://arxiv.org/abs/q-alg/9709040}{arXiv:q-alg/9709040} (1997).
本文では \S4.6 および \S6 の affine formality を用いる。

\bibitem{PDolgushev}
V.~Dolgushev,
The Van den Bergh duality and the modular symmetry of a Poisson variety,
\href{https://arxiv.org/abs/math/0612288v3}{arXiv:math/0612288v3} (2007).
Theorem~3 および Proposition~1。

\bibitem{Beilinson}
A. A. Beilinson,
\emph{Coherent sheaves on \(\PP^n\) and problems of linear algebra},
Functional Analysis and Its Applications \textbf{12} (1978), 214--216.
\url{https://doi.org/10.1007/BF01681436}.
原著Section 4の対角Koszul分解を確認。
\url{https://www.mathnet.ru/eng/faa2009}.
\bibitem{DDL}
O. De Deken and W. Lowen,
\emph{Abelian and derived deformations in the presence of \(\ZZ\)-generating geometric helices},
Journal of Noncommutative Geometry \textbf{5} (2011), 477--505.
Theorem 4.5; Propositions 5.5, 5.8, 5.9; Lemma 5.14.
\url{https://doi.org/10.4171/JNCG/83}.
\bibitem{Keller}
B. Keller, \emph{Deriving DG categories},
Annales scientifiques de l'École Normale Supérieure \textbf{27} (1994), 63--102.
\url{https://doi.org/10.24033/asens.1689}.
\bibitem{Mathew}
A. Mathew, \emph{The Galois group of a stable homotopy theory},
Advances in Mathematics \textbf{291} (2016), 403--541.
Proposition 3.28, Remark 3.29, Proposition 3.35.
\url{https://arxiv.org/abs/1404.2156}.
\bibitem{Orlov}
D. Orlov,
\emph{Derived categories of coherent sheaves and triangulated categories of singularities},
arXiv:math/0503632, v3 (2009), Theorem 2.5.
\url{https://arxiv.org/abs/math/0503632}.
\bibitem{CCKM}
T. Cassidy, A. Conner, E. Kirkman and W. F. Moore,
\emph{Periodic free resolutions from twisted matrix factorizations},
Journal of Algebra \textbf{455} (2016), 137--163.
Proposition 5.7, Theorem 5.8.
\url{https://doi.org/10.1016/j.jalgebra.2016.01.037}.
\bibitem{CT}
A. Căldăraru and J. Tu, \emph{Curved \(A_\infty\) algebras and Landau--Ginzburg models},
New York Journal of Mathematics \textbf{19} (2013), 305--342.
Section 5、特にdg Orlov windowのTheorem 5.8.
\url{https://nyjm.albany.edu/j/2013/19-16v.pdf}.
\bibitem{BFK}
M. Ballard, D. Favero and L. Katzarkov,
\emph{A category of kernels for equivariant factorizations and its implications for Hodge theory},
Publications Mathématiques de l'IHÉS \textbf{120} (2014), 1--111.
Theorem 5.39のKoszul cohomologyによる公式、Theorem 5.15とCorollary 5.18のequivariant外部積.
\url{https://doi.org/10.1007/s10240-013-0059-9}.
確認した公刊本文 \url{https://www.numdam.org/article/PMIHES_2014__120__1_0.pdf}.
\bibitem{LZ}
X. Lin and S. Zhang,
\emph{Serre algebra, matrix factorization and categorical Torelli theorem for hypersurfaces},
arXiv:2310.09927, v2 (2024), Theorem 1.7.
\url{https://arxiv.org/abs/2310.09927}.
\bibitem{HR}
D. Huybrechts and J. V. Rennemo,
\emph{Hochschild cohomology versus the Jacobian ring, and the Torelli theorem for cubic fourfolds},
Algebraic Geometry \textbf{6} (2019), 76--99.
Lemma 2.19およびSections 2--3.
\url{https://doi.org/10.14231/AG-2019-005}.
\bibitem{Huy}
D. Huybrechts, \emph{The K3 category of a cubic fourfold},
Compositio Mathematica \textbf{153} (2017), 586--620.
Proposition 3.12およびSection 3.5.
\url{https://doi.org/10.1112/S0010437X16008137}.
\bibitem{JLLZ}
A. Jacovskis, X. Lin, Z. Liu and S. Zhang,
\emph{Infinitesimal categorical Torelli theorems for Fano threefolds},
arXiv:2203.08187, Theorem 4.4(2).
\url{https://arxiv.org/abs/2203.08187}.
\bibitem{Lurie}
J. Lurie, \emph{Derived Algebraic Geometry X: Formal Moduli Problems},
2011. Theorems 5.1.9, 5.3.16およびProposition 5.3.21.
\url{https://www.math.ias.edu/~lurie/papers/DAG-X.pdf}.
\bibitem{BKP}
A. Blanc, L. Katzarkov and P. Pandit,
\emph{Generators in formal deformations of categories},
Compositio Mathematica \textbf{154} (2018).
本文で用いるのはarXiv版のDefinition 3.3およびSection 3.3.
\url{https://arxiv.org/abs/1705.00655}.
\bibitem{KL}
B. Keller and W. Lowen,
\emph{On Hochschild cohomology and Morita deformations},
International Mathematics Research Notices (2009).
\url{https://webusers.imj-prg.fr/~bernhard.keller/publ/KellerLowenHochDefo.pdf}.
\bibitem{EG}
P. Etingof and V. Ginzburg,
\emph{Noncommutative del Pezzo surfaces and Calabi--Yau algebras},
Journal of the European Mathematical Society \textbf{12} (2010), 1371--1416.
Sections 3.5--3.6, Theorem 3.6.2, Proposition 8.1.1(i) およびその証明。
\url{https://doi.org/10.4171/JEMS/235}.
\href{https://arxiv.org/abs/0709.3593v4}{arXiv:0709.3593v4}.
\bibitem{Walton}
C. Walton, \emph{Representation theory of three-dimensional Sklyanin algebras},
Nuclear Physics B \textbf{860} (2012), 167--185.
Theorem 2.14および式(2.15)--(2.16).
\url{https://doi.org/10.1016/j.nuclphysb.2012.02.015}.
\bibitem{CGWZ}
K. Chan, J. Gaddis, R. Won and J. J. Zhang,
\emph{Ozone groups of Artin--Schelter regular algebras satisfying a polynomial identity},
arXiv:2312.13014, v2 (2024). Introductionの非退化Sklyanin代数の記述。
\url{https://arxiv.org/abs/2312.13014}.
\bibitem{Elagin}
A. Elagin, \emph{On equivariant triangulated categories},
arXiv:1403.7027, Theorem 4.2.
\url{https://arxiv.org/abs/1403.7027}.
\bibitem{ELS}
A. Elagin, V. A. Lunts and O. M. Schn\"urer,
\emph{Smoothness of Derived Categories of Algebras},
Moscow Mathematical Journal \textbf{20} (2020), no. 2, 277--309.
確認した一次資料はarXiv:1810.07626v2 (22 March 2019),
Theorem 4.4 (Verdier局所化), Proposition 4.10とTheorem 4.15 (稠密開集合と生成子),
Theorem 4.18 (外部積生成), Theorem 5.1 (有限中心代数の有界導来圏のsmoothness).
\url{https://arxiv.org/abs/1810.07626};
\url{https://www.mathjournals.org/mmj/2020-020-002/2020-020-002-003.html}.

\bibitem{TV}
B. Toën and M. Vaquié,
\emph{Moduli of objects in dg-categories},
Annales scientifiques de l'École Normale Supérieure (4) \textbf{40} (2007), 387--444.
Theorem 3.6; Corollaries 2.13, 3.15, 3.19, 3.21; Lemma 3.25 and proof of Corollary 3.24.
\url{https://arxiv.org/abs/math/0503269};
確認した著者版 \url{https://www.math.univ-toulouse.fr/~btoen/ttt.pdf}.
\bibitem{Stacks}
The Stacks Project Authors, \emph{The Stacks Project}.
Tag 02HX (Lemma 37.12.1, Noetherian底上の無限小smoothness判定),
Tag 055U (Lemma 37.38.6, smooth morphismのétale局所section).
\url{https://stacks.math.columbia.edu/tag/02HX};
\url{https://stacks.math.columbia.edu/tag/055U}.
参照日2026年11月4日。
\bibitem{Pirozhkov}
D. Pirozhkov, \emph{A categorical Torelli theorem for hypersurfaces},
Bulletin of the London Mathematical Society \textbf{56} (2024), 3075--3089.
\url{https://doi.org/10.1112/blms.13117}.
公刊版のabstractで、cubicに対する仮定 \(n>3\) を確認。
本稿の全次元の一意性には \cite[Theorem 1.7]{LZ} を用いる。
\bibitem{HuyUpdate}
D. Huybrechts, \emph{The K3 category of a cubic fourfold: an update},
Beiträge zur Algebra und Geometrie / Contributions to Algebra and Geometry (2025).
Section 9: Unknown deformations.
\url{https://doi.org/10.1007/s13366-025-00788-9}.

\bibitem{ChanMulti}
D. Chan, \emph{Twisted multi-homogeneous coordinate rings},
Journal of Algebra \textbf{223} (2000), 438--456.
Theorem 3.4; proof of Theorem 5.2, Claim 5.3 and Lemmas 5.4--5.6;
Corollaries 5.8--5.9.
\url{https://doi.org/10.1006/jabr.1999.7938}.
確認した著者版 \url{https://web.maths.unsw.edu.au/~danielch/paper/tbcr.ps}, pp.8--9,14--17.
\bibitem{SmithTate}
S. P. Smith and J. Tate,
\emph{The center of the 3-dimensional and 4-dimensional Sklyanin algebras},
K-Theory \textbf{8} (1994), 19--63.
Corollary 2.3 (p.~26); Theorem 3.7(c) (p.~38).
\url{https://doi.org/10.1007/BF00962089};
\url{https://sites.math.washington.edu/~smith/Research/tate.pdf}.
\bibitem{KeelerMulti}
D. S. Keeler, \emph{Noncommutative ampleness for multiple divisors},
Journal of Algebra \textbf{265} (2003), no. 1, 299--311.
Theorem 2.7, Corollary 2.8 (任意有限個のbimoduleのNC-ampleness);
Hypothesis 3.1, Theorem 3.3, Corollary 3.5 (二因子のnoetherian性).
番号はarXiv一次TeXソースで確認し、刊行版との番号一致は未照合。
\url{https://doi.org/10.1016/S0021-8693(03)00126-1};
\url{https://arxiv.org/abs/math/0210417};
\url{https://arxiv.org/e-print/math/0210417}.
\bibitem{WWY3}
C. Walton, X. Wang and M. Yakimov,
\emph{Poisson geometry of PI three-dimensional Sklyanin algebras},
Proceedings of the London Mathematical Society \textbf{118} (2019), 1471--1500.
Theorem 1.3(1),(2a), verified in arXiv:1704.04975v3, p.~3.
\url{https://doi.org/10.1112/plms.12220};
\url{https://arxiv.org/abs/1704.04975}.
\bibitem{ToenFinitude}
B. Toën, \emph{Finitude homotopique des dg-algèbres propres et lisses},
Proceedings of the London Mathematical Society (3) \textbf{98} (2009), 217--240.
Lemmas 3.13--3.14. Lemma 3.14の結論は部分前層のquasi-compactnessであり、一般の開性ではない。
\url{https://doi.org/10.1112/plms/pdn027};
確認した著者版 \url{https://www.math-beb.fr/wp-content/uploads/2026/10/finitesatdgalg.pdf}.

\bibitem{SnowdenDVR}
A. Snowden, \emph{Lecture 8: Elliptic curves over DVRs},
Math 679 lecture notes, University of Michigan, 2013.
Section ``Reduction of torsion points'', good reductionでのprime-to-residue-characteristic torsionの還元の単射性。
\url{https://websites.umich.edu/~asnowden/teaching/2013/679/L08.html}.
参照日2026年11月6日。
\bibitem{HMY}
J.-W. He, X.-C. Ma and Y. Ye, \emph{Generalized Kn\"orrer's Periodicity Theorem},
arXiv:2107.06438v2 (16 August 2021).
対象は中心二次元の商であり、本稿の三次の場合には直接適用しない。
\url{https://arxiv.org/abs/2107.06438}.

% Current verified additions, 2026-11-08.
% Insert inside the manuscript's thebibliography environment.

\bibitem{UeyamaIsolated}
K. Ueyama,
\emph{Noncommutative graded Gorenstein isolated singularities},
in \emph{Proceedings of the 45th Symposium on Ring Theory and Representation Theory}
(Shinshu University, September 7--9, 2012),
K. Koike (ed.), February 2013, 158--163.
Theorem 7, printed p.~161.
\url{https://ring-theory-japan.com/ring/oldmeeting/2012/report2012/proceedings.pdf}.

\bibitem{ZhangTwist}
J.~J. Zhang,
\emph{Twisted graded algebras and equivalences of graded categories},
Proceedings of the London Mathematical Society (3) \textbf{72} (1996), no.~2, 281--311.
Theorems 1.1 and 1.3.
\url{https://doi.org/10.1112/plms/s3-72.2.281}.
\url{https://deepblue.lib.umich.edu/bitstream/handle/2027.42/135651/plms0281.pdf}.

\bibitem{RRZ}
M. Reyes, D. Rogalski and J.~J. Zhang,
\emph{Skew Calabi--Yau algebras and homological identities},
Advances in Mathematics \textbf{264} (2014), 308--354.
Theorems 0.3 and 5.4; Example 5.5.
\url{https://doi.org/10.1016/j.aim.2014.07.010}.
\url{https://arxiv.org/abs/1302.0437}.

\bibitem{Griffiths}
P.~A. Griffiths,
\emph{On the periods of certain rational integrals.~I},
Annals of Mathematics (2) \textbf{90} (1969), no.~3, 460--495.
Section 8, especially Theorems 8.1 and 8.3 and formula (8.6).
\url{https://doi.org/10.2307/1970746}.
\url{https://www.cimat.mx/~luis/seminarios/Teoria-de-Hodge/period_integral_I.pdf}.

\bibitem{KuznetsovFCY}
A. Kuznetsov,
\emph{Calabi--Yau and fractional Calabi--Yau categories},
Journal f\"ur die reine und angewandte Mathematik \textbf{753} (2019), 239--267.
Lemmas 2.6, 3.13 and 3.19; Corollary 4.1.
\url{https://doi.org/10.1515/crelle-2017-0004}.

\bibitem{CHL}
C.-H. Cho, H. Hong and S.-C. Lau,
\emph{Noncommutative homological mirror functor},
Memoirs of the American Mathematical Society \textbf{271} (2021), no.~1326.
Chapter 5, especially Theorems 5.15 and 5.16.
\url{https://doi.org/10.1090/memo/1326}.
\url{https://arxiv.org/abs/1512.07128}.

\bibitem{Lee}
S. Lee,
\emph{Noncommutative homological mirror symmetry of elliptic curves},
Kyoto Journal of Mathematics \textbf{61} (2021), no.~3, 723--743.
Theorem 6.3 and the corollary in the introduction.
\url{https://arxiv.org/abs/2004.09145}.

\bibitem{CisinskiTabuada}
D.-C. Cisinski and G. Tabuada,
\emph{Lefschetz and Hirzebruch--Riemann--Roch formulas via noncommutative motives},
Journal of Noncommutative Geometry \textbf{8} (2014), 1171--1190.
Section 4 and Example 5.6.
\url{https://doi.org/10.4171/JNCG/183}.
\url{https://www.math.univ-toulouse.fr/~dcisinsk/ncLHRR.pdf}.

\bibitem{MathewHH}
A. Mathew,
\emph{Kaledin's degeneration theorem and topological Hochschild homology},
Geometry \& Topology \textbf{24} (2020), no.~6, 2675--2708.
Theorem 1.3.
\url{https://doi.org/10.2140/gt.2020.24.2675}.
\url{https://arxiv.org/abs/1710.09045}.

\bibitem{PerryIHC}
A. Perry,
\emph{The integral Hodge conjecture for two-dimensional Calabi--Yau categories},
Compositio Mathematica \textbf{158} (2022), no.~2, 287--333.
Lemmas 3.4 and 3.6; Theorem 3.5.
\url{https://doi.org/10.1112/S0010437X22007266}.

\end{thebibliography}

\endgroup
\end{document}
